\providecommand\ExamFontPath{../../../字体/}
\documentclass[fontset=exam]{exam-zh}
\usepackage{longtable,booktabs,array,calc,multirow,fvextra}
\usepackage[normalem]{ulem}
\providecommand{\hl}[1]{#1}
\providecommand{\ul}[1]{\uline{#1}}
\DefineVerbatimEnvironment{verbatim}{Verbatim}{breaklines=true,fontsize=\small}
\providecommand{\tightlist}{\setlength{\itemsep}{0pt}\setlength{\parskip}{0pt}}
\newcommand{\pandocbounded}[1]{#1}
\setlength{\emergencystretch}{3em}
\setlength{\tabcolsep}{2pt}
\examsetup{page/size=a4paper}
\title{厦门大学 2011—2012 年 期中考试}
\subject[4em]{热学}
\begin{document}
\maketitle
\fancyfoot[C]{第 \thepage 页（共 \pageref*{LastPage} 页）}
1．（15 分）

一氧气瓶容积为 3.2×10\textsuperscript{-2}m\textsuperscript{3}，瓶内氧气压强为 130atm，按规定瓶内氧气压强降至 10atm 时就得重新充气，某单位每天消耗 1atm 的氧气 0.4m\textsuperscript{3}，设氧气使用过程中温度保持不变，问这瓶氧气在该单位能用几天？

2．（20 分）

容器中装有 4mol 标态下的氧气。求：

（1）容器中单位体积的分子数；

（2）每个分子的平均平动动能；

（3）分子的平均速率；

（4）系统的内能；

（5）分子运动的平均自由程。

（设氧分子的有效直径\(d = 3.6 \times 10^{- 10}m\)）

3．（15 分）

储有氧气的容器以速度\(v = 100m/s\)运动，氧气的压强\(p = 1atm\)，温度为\(T = 300K\)。该容器突然停止，全部定向运动的动能变成气体分子热运动的动能。求：

（1）系统的温度改变了多少？

（2）系统的压强改变了多少？

（设容器的体积保持不变）

4．（18 分）

某一假想的气体系统含有\emph{N}个气体分子，单个分子质量为\emph{m}，其分子速率分布函数为：

\[f(v) = \left\{ \begin{matrix}
{C(v - v_{0})(2v_{0} - v)\quad(v_{0} \leq v \leq 2v_{0})} \\
{0\quad\quad\quad\;\;(v < v_{0}\text{或}v > 2v_{0})}
\end{matrix} \right.\]

其中\(v_{0}\)为已知量，\emph{C}为归一化常数。求：

（1）\emph{C}=?；

（2）速率在\(\left. 0 \right.\sim(5/4)v_{0}\)之间的分子数；

（3）速率在\(\left. 0 \right.\sim(5/4)v_{0}\)之间的分子的平均平动动能。

5．（20 分）

分子被约束在 xy 平面内运动，已知分子速度在\(\left. v_{x} \right.\sim v_{x} + \text{d}v_{x}\)、\(\left. v_{y} \right.\sim v_{y} + \text{d}v_{y}\)内的分子数占总分子数的百分比为

\[\frac{dN_{\overset{\rightarrow}{v}}}{N} = \frac{m}{2\pi kT}e^{- \frac{m(v_{x}^{2} + v_{y}^{2})}{2kT}}dv_{x}dv_{y}\]

式中\(m\)为分子质量，\(T\)为系统温度，\(k\)为玻尔兹曼常数。试求：

\begin{enumerate}
\def\labelenumi{\arabic{enumi}.}
\item
  分子速率在\(\left. v \right.\sim v + \text{d}v\)内的分子数占总分子数的百分比；
\item
  分子平均平动动能（\(\varepsilon = \frac{1}{2}mv^{2}\)）在\(\left. \varepsilon \right.\sim\varepsilon + \text{d}\varepsilon\)内的分子数占总分子数的百分比；
\item
  分子的平均平动动能。
\end{enumerate}

6．（12 分）

一柱型容器底面积为\emph{S}，高度为\emph{H}，内含有\emph{N}个气体分子，每个分子的质量为\emph{m}，系统的温度为\(T\)，该容器沿竖直方向放置。考虑到重力的作用，且重力加速度为\emph{g}，试分别求出容器底部和顶部的分子数密度。

\end{document}
